
\begin{frame}{Point de départ: relation globale et dépendances}

On part d'un schéma contenant tous les attributs connus.

\vfill
\centerline{\texttt{(code, nom, adresse, activité, description)}}
\vfill

On identifie les dépendances fonctionnelles

\begin{redbullet}
  \item $code \to nom, adresse$
  \item $adresse \to code, nom$
 \item $code, activit\acute{e}  \to  description$
\end{redbullet}

\vfill

On va baser la décomposition sur les dépendances fonctionnelles.

\end{frame}


\begin{frame}{Décomposition guidée par les dépendances}

Trois DFs, trois ensembles d'attributs
\,: (\textbf{code}, nom, adresse), (code, nom, \textbf{adresse}), (\textbf{code, activité}, description)

\vfill

\begin{redbullet}
  \item  Les deux premiers définissent une relation 
  
\smallskip  
      \centerline{Logement (\textbf{code}, nom, adresse)}
\smallskip
   Deux clés dites \alert{candidates},   on en choisit une comme \alert{clé primaire}
 \item Le second définit une relation 
 
\medskip

 \centerline{Activité (\textbf{code, activité}, description)}
 
\end{redbullet}

\vfill
Relations en 3FN, décomposition \alert{sans perte d'information}



\end{frame}

%%%%%%%%%%%%
\begin{frame}{Regardons les séjours}

Le schéma global de départ est le suivant:


\vfill
\begin{small}
\texttt{(code, nom, adresse, idVoyageur, nomV, prénomV)}
\end{small}
\vfill

On a les dépendances suivantes :
\begin{redbullet}
  \item $code \to nomL$
  \item $idVoyageur \to pr\acute{e}nomV, nomV$
\end{redbullet}

\vfill

La clé est le couple $(code, idVoyageur)$. 

\vfill
\begin{remblock}{En troisième forme normale?}
Non, car pour les \alert{deux}  DF,  la partie gauche n'est pas la clé
\end{remblock}
\end{frame}

%%%%%%%%%%%%
\begin{frame}{La décomposition}

Comme avant, à partir des DF.

\vfill


\begin{redbullet}
  \item On retrouve Logement (\textbf{code}, nom, adresse) 
  \item Voyageur (\textbf{idVoyageur}, nom, prénom)
\end{redbullet}

\vfill

\alert{Pas suffisant}, on a perdu le lien entre les logements et les voyageurs.

\vfill
On ajoute une relation avec la clé.

\begin{redbullet}
  \item Séjour (\textbf{idVoyageur, code})
\end{redbullet}

\vfill
\alert{On est en 3FN, sans perte d'information (jointures pour reconstituer)}

\end{frame}

%%%%%%%%%%%%
\begin{frame}{Algorithme de normalisation}

On part d’un schéma de relation $R$ \alert{global}
et d’un ensemble de dépendances fonctionnelles \alert{minimales} et \alert{directes} (exercice)

\vfill

On détermine alors les clés de R

\vfill 
\begin{redbullet}
  \item Pour chaque DF $X \to A_1 \cdots, A_n$,
      on crée une relation $(X,A_1 \cdots, A_n)$ de clé $X$
     \item On unifie les relations de même clé et les relations de même schéma
  \item Pour chaque clé $C$
non représentée dans une des relations précédentes, on crée une relation $(C)$
\end{redbullet}

\vfill

\alert{On obtient un schéma normalisé}
\end{frame}


%%%%%%%%%%%%
\begin{frame}{Et en pratique ?}

Dans la vraie vie, les DF n'existent pas naturellement. On a plutôt\,:

\vfill
\begin{small}
\texttt{(nomLogement, adresse, nomVoyageur, typeLogement, ville)}
\end{small}
\vfill 

Pas de DF.... Une phase d'analyse détermine 
des \alert{entités}, leurs identifiants et 
leurs liens.

\vfill

Ici, entités Logement et Voyageur, avec\,:

\begin{redbullet}
  \item $code \to nomLogement, adresse, typeLogement$
    \item $idVoyageur \to nomVoyageur, ville, pays$
  
\end{redbullet}

Maintenant on normalise et on obtient un schéma en 3FN.
\end{frame}

%%%%%%%%%%%%%
\begin{frame}{À retenir}

Il est toujours possible de se ramener à un schéma normalisé.

\vfill
\begin{redbullet}
   \item En déterminant les DF et les clés
   \item En appliquant l'algorithme de normalisation
\end{redbullet}

\vfill

\alert{En pratique}: les DF ne sont pas données naturellement. Il faut les déterminer
par un processus d'analyse des besoins et de  conception.


\end{frame}

